\section{Conclusion}
\label{sec:conclusion}

A three-dimensional thermo-mechanical finite element workflow for ship-plate line heating has been presented and assessed against an available subset of Li et al.\ \cite{li2023} Q345 benchmark data. The workflow includes transient heat conduction with moving Gaussian and reduced full-line thermal-flux source options, temperature-dependent Q345 properties (Table~1), convection, quenching, optional radiation, linear thermoelasticity, and a calibrated inherent-strain surrogate for residual bending. The validated quantity is the final residual-deflection response of this reduced workflow; the paper does not claim full transient elastoplastic residual-stress validation.

Validation on eleven available Table~2 experiments (cases 1--7 and 87--90) yields MAPE \SI{1.17}{\%} in midspan deflection (maximum \SI{2.78}{\%}) for the moving-torch baseline; the four-case calibration subset achieves MAPE \SI{0.52}{\%}, while the seven non-anchor validation cases achieve MAPE \SI{1.55}{\%}. Multi-pass and high-energy conditions are predicted within \SI{0.8}{\%}. An exploratory velocity-dependent post-calibration using Cases~1--7 reduces the speed-sweep MAPE from \SI{1.67}{\%} to \SI{0.71}{\%}, but this correction is not treated as an independently validated model. The full-line thermal-flux study demonstrates that pass scheduling is essential: a collapsed quench-enabled full-line heat underpredicts two-pass deflections by about \SI{50}{\%}, while sequential line heating with quench matches the multi-pass subset with MAPE \SI{0.27}{\%}. Keel-plate cases and the skin-depth thermal source are numerical process-planning demonstrations only; they are not claimed as experimental validation.

The previously identified verification items have been advanced where the archived dataset permits: the velocity-dependent inherent-strain calibration, keel-plate thermal-unit audit, mesh/energy reporting scaffold, skin-depth thermal source, a Case~6 mesh/time-step sensitivity matrix, and a diagnostic Case~6 J2 elastoplastic spot check are generated reproducibly from repository outputs. The mesh/time-step study shows that local heating-band mesh resolution dominates the Case~6 discretization sensitivity, while halving $\Delta t$ at the baseline mesh changes $|w|_{\max}$ by only \SI{0.03}{\%}. A more complete elastoplastic validation would still require further development because the present J2 spot check does not reproduce the calibrated residual-deflection magnitude. Measured keel-plate validation and a full electromagnetic coil-field model require new experiments or physics modules. Table~2 rows 8--86 remain unavailable in the current source excerpt. All inputs and outputs required to reproduce the completed results are archived in the accompanying project repository.
